%------------------------------------------------------------------------------%
\section{Exercícios}

\begin{exercises*}

%------------------------------------------------------------------------------%
\begin{exercise}
  Faça o esboço gráfico das funções abaixo
  \begin{mathlist}
    \item f(x) = 2^x-2
    \item f(x) = \left( \frac{1}{5} \right)^x
    \item f(x) =-4\left( \frac{1}{2} \right)^x
    \item f(x) = 3^{-x}
  \end{mathlist}

  \begin{answer}
    \begin{alphalist}
        \item \IncludeGraphicCentred{gabarito_aula4-1a}[0.9]
        \item \IncludeGraphicCentred{gabarito_aula4-1b}[0.9]
        \item \IncludeGraphicCentred{gabarito_aula4-1c}[0.9]
        \item \IncludeGraphicCentred{gabarito_aula4-1d}[0.9]
    \end{alphalist}
  \end{answer}
\end{exercise}

%------------------------------------------------------------------------------%
\begin{exercise}
  Encontre o domínio das seguintes funções:
  \begin{mathlist}
    \item g(x) = \frac{-5}{ \sqrt{16-\left( \frac{1}{2} \right)^x}}
    \item f(x) = \sqrt{9 \cdot 3^{2x-1}-\frac{1}{27}}.
  \end{mathlist}

  \begin{answer}
    \begin{mathlist}
      \item \{ x \in \R \st x    > -4 \}
      \item \{ x \in \R \st x \geq -2 \}
    \end{mathlist}
  \end{answer}
\end{exercise}

%------------------------------------------------------------------------------%
\begin{exercise}
  Resolva as equações abaixo
  \begin{mathlist}
    \item 8^{x} = 0,25
    \item 5^{2x-7} = 125
    \item 8^{x^2-x} = 4^{x+1}
    \item 2^{x^2- 3x - 4} = 1
    \item \sqrt{\left(\frac{1}{2}\right)^{3x-2}}
           = \left(\frac{1}{2}\right)^{-4x}.2^{-x+4}
    \item (10^x)^{1-x}-0,000001 = 0
    \item 2^{2x + 1} - 3.2^{x + 2} = 32
    \item 2^{x^2+5}+2^4 = 144
    \item 5^{4x-1}-5^{4x}-5^{4x+1}+5^{4x+2} = 480
    \item 2^{3x-2}-4^{x+6} = 0
  \end{mathlist}

  \begin{answer}
    \begin{mathlist}
      \item S = \left\{ -\frac{2}{3}        \right\}
      \item S = \left\{ 5                   \right\}
      \item S = \left\{ -\frac{1}{3}, 2     \right\}
      \item S = \left\{ -1, 4               \right\}
      \item S = \left\{ -\frac{2}{3}        \right\}
      \item S = \left\{ -2, 3               \right\}
      \item S = \left\{ 3                   \right\}
      \item S = \left\{ -\sqrt{2}, \sqrt{2} \right\}
      \item S = \left\{ \frac{1}{2}         \right\}
      \item S = \left\{ 14                  \right\}
    \end{mathlist}
  \end{answer}
\end{exercise}

%------------------------------------------------------------------------------%
\begin{exercise}
  Resolva as inequações abaixo
  \begin{mathlist}
    \item \left(\frac{1}{2}\right)^{2x-1}\leq \frac{1}{16}
    \item e^{x^2-5}\geq e^{4x}
    \item \left(\frac{1}{2} \right)^{2x-1}> \frac{1}{16}
    \item 7^{x-2}<343
    \item 1 \leq \left( \frac{1}{2}\right)^x \leq 8
  \end{mathlist}

  \begin{answer}
    \begin{mathlist}
      \item \left\{ x \in \R \st x \geq \frac{5}{2}            \right\}
      \item \left\{ x \in \R \st x \leq -1 \text{ ou } x\geq 5 \right\}
      \item \left\{ x \in \R \st x < \frac{5}{2}               \right\}
      \item \left\{ x \in \R \st x < 5                         \right\}
      \item \left\{ x \in \R \st -3 \leq x \leq 0              \right\}
    \end{mathlist}
  \end{answer}
\end{exercise}

%------------------------------------------------------------------------------%
\begin{exercise}
  Se \(f(x) = 16^{1+\sfrac{1}{x}}\),
  encontre o valor de $f(-1) + f(-2) + f(-4)$

  \begin{answer}
    13
  \end{answer}
\end{exercise}

%------------------------------------------------------------------------------%
\begin{exercise}
  Calcule $f(0)-f (3/2)$ sabendo que
  \[
    f(x) = \begin{cases}
      2^x,         & \text{ se } -1 \leq x \leq 1, \\
      \frac{1}{x}, & \text{ se } x > 1.
    \end{cases}
  \]

  \begin{answer}
    \(\frac{1}{3}\)
  \end{answer}
\end{exercise}

%------------------------------------------------------------------------------%
\begin{exercise}
  Seja a função $f\colon \R\to\R$ definida por $f(x) = 2^x$.
  Encontre o valor de $f(a+1)-f(a)$

  \begin{answer}
    $2^a$
  \end{answer}
\end{exercise}

\end{exercises*}

%------------------------------------------------------------------------------%